% s-194-p194-m002; original PDF 194, 195
\Aphorism{s-194-p194-m002}{281}
Wie? die Stoffwahl des Künstlers wäre gleichgültig? Gleichgültig, ob ein Dichter an historische Opernhelden oder an grossstädtische Mischmaschgewächse oder an Menschen grosser Leidenschaft, tiefer Erkenntniss, gewaltigen Schicksals seine Darstellungskunst wendet — ob es sich um Pescaras oder einer Näherin Versuchung handelt, ob der Fall Tristan und Isolde oder der Fall Hans und Grete vorliegt? Nur auf das Wie, nicht auf das Was käme es an? — Man muss sich erinnern, dass zwischen Künstler und Kunstfreund, zwischen Geber und Empfänger eine seltsame Art Mensch als Vermittler und nicht immer ehrlicher Makler ihr Wesen treibt; auch die Kunst hat ihre Börse, die Ästhetik, wo Leute auf dem Papier Getreide kaufen und verkaufen, die nie ein wirkliches Weizenkorn gesehen haben. Die Ästhetiker, will sagen Menschen, an deren künstlerische Unfruchtbarkeit und Unempfänglichkeit man bis zum Beweise des Gegentheils glauben soll, haben jene sinnlose Trennung zwischen dem Was und Wie der künstlerischen Darstellung auf den Markt philosophischer Begriffe gebracht, zum grössten Erstaunen der eigentlichen Producenten und Consumenten. Kein Zuschauer bringt es fertig, die Form „Hamlet“ ohne den Inhalt, oder den Inhalt „Hamlet“ ohne die Form für sich zu geniessen, etwa wie man die Schale von der Frucht trennen und „abstrahiren“ kann. Kein Künstler hat sich von jener Willensfreiheit etwas träumen lassen, die es ihm anheimgäbe, aus der Urne aller vorhandenen Kunstformen und aus der aller vorhandenen Stoffe je ein Exemplar herauszuziehen und beide Bestandtheile zu einem Kunstwerke zusammenzukitten. Ihm wäre mit dieser Freiheit, die gleich ihrer moralischen Zwillingsschwester die Vorstellung Blinder von der Farbe ist, nicht einmal gedient; er würde auch unter tausend historisch überlieferten Stoffen und Formen nicht finden, was er zu seinem tausend und ersten Kunstwerk braucht, und darum zieht er es vor, mit Umgehung dieses ganzen Wahlgeschäfts, das er den Epigonen und ewig zu spät Kommenden überlässt, schlankweg sein Erlebniss in seiner Sprache zu verkünden. Damit ist nicht eine neue Form und ein neuer Inhalt, sondern eine neue untrennbare Einheit von Form und Inhalt, eine neue organische Durchdringung des Was mit dem Wie geschaffen und den Ästhetikern eine neue Gelegenheit gegeben, ein Ganzes entzwei zu denken.
\par

